\documentclass[11pt]{article}
\usepackage[margin=1.1in]{geometry}
\usepackage{amsmath, amssymb, amsthm}
\usepackage{booktabs}
\usepackage{listings}
\usepackage{xcolor}
\usepackage{hyperref}
\hypersetup{colorlinks=true, linkcolor=blue!50!black, urlcolor=blue!50!black}

\lstset{
  basicstyle=\ttfamily\footnotesize,
  breaklines=true,
  frame=leftline,
  framesep=6pt,
  xleftmargin=10pt,
  rulecolor=\color{gray!50},
  backgroundcolor=\color{gray!5},
}

\newcommand{\code}[1]{\texttt{#1}}
\DeclareMathOperator{\hn}{HN}
\DeclareMathOperator{\Var}{Var}
\DeclareMathOperator{\sd}{sd}
\DeclareMathOperator{\E}{\mathbb{E}}

\title{\textbf{A stale comment, a silent failure}\\[2pt]
\large Why \code{z.init = 0.01} destroyed the time-block forecast study}
\author{Technical note}
\date{12 July 2026 \quad(rev.\ 18 July 2026)}

\begin{document}
\maketitle

\begin{abstract}
The time-block forecast study
(\code{code/analysis/time\_block\_forecast/simstudy/}) reported that the AR(2)
model forecasts \emph{four times worse} than an i.i.d.-in-time baseline, with
the gap growing monotonically in lead time. That result is an artifact of a
single argument, \code{z.init = 0.01}, passed to \code{mcmc()} on the strength
of a comment that was true when it was written and false by the time it was
read. This note explains why the comment was wrong, derives the cascade from
that one number to a predictive standard deviation of $104$ against a data
marginal of $4.38$, and documents the verification --- including the effect on
the study's own scores, where the fix halves the AR(2) CRPS and cuts its Brier
by a third. It then records two findings from the cleanup: the model
non-identifiability the bug exploited persists as a milder, stochastic
``reflected'' branch that corrupts \emph{every} proper score (Brier included,
not only the multivariate ones), and the same fatal default still lies dormant
behind a guard in a sibling backend. The deeper point is not the bug: it is
that the workaround converted a \emph{loud} failure into a \emph{silent} one.
\end{abstract}

\section{The comment}

\code{run-settings.R} carried this justification for overriding the default
initial value of the latent skew process $z$:

\begin{lstlisting}
# z.init = 0 (the backend default) is fatal for the temporal/AR(2) z
# update: hn.cop(0, sig) = qnorm(phn(0, sig)) = -Inf, the MH ratio R
# becomes NaN, and the `if (!is.na(R))` guard in updateZTS_AR2
# silently rejects every proposal -- z and phi.z stay stuck at 0,
# producing all-NaN forecasts. A small positive initial value lets
# hn.cop evaluate finitely on the very first iteration.
fit <- mcmc(..., z.init = 0.01, ...)
\end{lstlisting}

Every sentence about the \emph{mechanism} is correct. The sentence about the
\emph{default} is not.

\section{What the backend actually does}

The model couples $z$ to a latent Gaussian process through a half-normal
copula. With $\sigma = \tau^{-1/2}$,
\begin{equation}
  z^\star \;=\; \operatorname{hn.cop}(z, \sigma)
        \;=\; \Phi^{-1}\!\big(\operatorname{phn}(z,\sigma)\big)
        \;=\; \Phi^{-1}\!\big(2\,\Phi(z/\sigma) - 1\big).
  \label{eq:hncop}
\end{equation}
At $z = 0$ this gives $\Phi^{-1}(0) = -\infty$, and the Metropolis ratio
degenerates to \code{NaN}. That failure was real, and it was \emph{fixed in the
backend}. From \code{code/R/ar2/mcmc\_ar2.R:249--250}:

\begin{lstlisting}
if (is.null(z.init)) {
    z.init <- 0.6745 / sqrt(tau.init)  # median of HN(1/sqrt(tau.init)) -> z.star = 0
}
\end{lstlisting}

The signature default is \code{z.init = NULL}, \emph{not} $0$. When \code{NULL},
$z$ is initialised at the \textbf{median of its own half-normal marginal}. This
choice is exact: since $0.6745 = \Phi^{-1}(0.75)$, substituting
$z = 0.6745\,\sigma$ into \eqref{eq:hncop} gives
\begin{equation}
  z^\star \;=\; \Phi^{-1}\!\big(2\Phi(0.6745) - 1\big)
          \;=\; \Phi^{-1}(2 \times 0.75 - 1)
          \;=\; \Phi^{-1}(0.5) \;=\; 0 ,
\end{equation}
\emph{dead centre} of the copula's Gaussian scale, for any $\tau$. The backend's
own \code{README.md} documents this fix, its symptom, and its rationale --- and
closes with the sentence that seals the trap:

\begin{quote}
\itshape ``Callers who pass \code{z.init} explicitly are unaffected.''
\end{quote}

\section{Why the comment was true once, and false later}

The chronology is the whole story.

\begin{enumerate}\itemsep2pt
  \item The backend default \emph{was} $z.init = 0$. It produced $-\infty$, then
        \code{NaN}, then a frozen chain. This is the bug the comment describes.
  \item \code{run-settings.R} worked around it locally by passing
        \code{z.init = 0.01}. The $-\infty$ disappeared; the study ran.
  \item The backend was then fixed \emph{properly}, changing the default to
        \code{NULL} $\mapsto$ the half-normal median, and documenting it.
  \item The correct fix could not reach \code{run-settings.R}, precisely because
        it passed \code{z.init} explicitly --- exactly the case the README
        excludes.
  \item Nobody removed the now-obsolete workaround, and its comment --- an
        accurate description of a world that no longer existed --- kept the
        workaround looking justified.
\end{enumerate}

The comment did not merely go stale. It \emph{actively defended} the line that
had to be deleted.

\section{What $0.01$ actually does}

Take \code{tau.init = 1}, so $\sigma = 1$. Equation~\eqref{eq:hncop} gives

\begin{center}
\begin{tabular}{lcl}
\toprule
\code{z.init} & $z^\star = \operatorname{hn.cop}(z,1)$ & \\
\midrule
$0$                      & $-\infty$ & the original bug: \code{NaN}, chain frozen \\
$0.01$ \emph{(the workaround)} & $\mathbf{-2.41}$ & finite, but a $2.4\sigma$ tail start \\
$0.6745$ \emph{(the correct default)} & $\mathbf{0}$ & the centre of the copula scale \\
\bottomrule
\end{tabular}
\end{center}

The workaround did not restore the chain to a sensible starting point. It moved
it from an \emph{impossible} one to an \emph{extremely improbable} one. A
random-walk Metropolis chain started $2.4$ standard deviations into the tail of
the very space it must explore does not, in practice, come back.

\section{The cascade}

\subsection{Step 1: $z$ freezes}

Started at $z^\star = -2.41$, the chain never reaches the bulk. Across $20{,}000$
iterations the fitted $z$ stays at $\approx 0.2$, while its own model
distribution, $z \sim \hn(0,\sigma)$ with $\sigma \approx 2$, demands
\begin{equation}
  \E[z] \;=\; \sigma\sqrt{2/\pi} \;\approx\; 1.6 .
  \label{eq:zmean}
\end{equation}
The fitted $z$ is a factor of $7$--$8$ below the value the model itself
prescribes. Its maximum over the entire retained chain is $0.30$; in a healthy
fit it reaches $15$.

\subsection{Step 2: the AR coefficient chases the freeze}

A frozen $z$ is, as a time series, nearly constant --- and a nearly constant
series looks maximally persistent. The AR(2) update on $z^\star$ obliges,
driving $\hat\phi_{1z} \to 1.11$ (truth: $0.80$), i.e.\ toward a unit root. The
near-unit-root prior then pins $z$ \emph{harder}. The freeze is
self-reinforcing: $\phi \mid z \to$ unit root if $z$ looks smooth; $z \mid \phi
\to$ smooth if $\phi$ is a unit root.

\subsection{Step 3: $\lambda$ and $\beta_0$ compensate}

The data see the skew channel only through the \emph{product} $\lambda z$, in
\begin{equation}
  \mu \;=\; \beta_0 \;+\; \lambda\,\bar z .
  \label{eq:ridge}
\end{equation}
This is a ridge: $(\beta_0, \lambda, \bar z)$ can slide along it freely provided
$\mu$ is preserved. With $\bar z$ pinned $7\times$ too small, the sampler
restores $\mu$ by sending $\lambda$ and $\beta_0$ to absurd values. All three
blocks reproduce the data mean of $14.5$ \emph{exactly}:

\begin{center}
\begin{tabular}{lrrrl}
\toprule
 & $\beta_0$ & $\lambda$ & $\bar z$ & $\mu = \beta_0 + \lambda\bar z$ \\
\midrule
truth                  & $10$    & $+3$    & $\approx 1.6$ & $\approx 14.5$ \\
block 1 (healthy)      & $10.36$ & $+2.91$ & $1.42$ & $\mathbf{14.50}$ \\
block 3 (\emph{broken}) & $36.41$ & $\mathbf{-97.75}$ & $0.22$ & $\mathbf{14.92}$ \\
\bottomrule
\end{tabular}
\end{center}

\textbf{This is why the failure is silent.} The likelihood is satisfied. The
fitted mean is right. Nothing inside the fit looks wrong. Only the individual
parameters are catastrophic --- and $\lambda$ has even flipped sign.

\subsection{Step 4: the forecast exposes the inconsistency}

\code{forecast\_block} is correct, and that is exactly why it detonates. It does
not reuse the frozen $z$; it redraws $z$ from the model's own law,
$z_{\text{fc}} \sim \hn(0, \tau_{\text{fc}}^{-1/2})$, which by
\eqref{eq:zmean} has mean $\approx 1.6$. It then multiplies by the corrupted
$\lambda$:
\begin{equation}
  \lambda \, z_{\text{fc}} \;\approx\; (-97.75)(1.6) \;\approx\; -156 ,
  \qquad
  \operatorname{sd}\big(\lambda z_{\text{fc}}\big)
    \;=\; |\lambda|\,\sigma\sqrt{1 - 2/\pi} \;\approx\; 118 .
\end{equation}
Against a data marginal standard deviation of $4.38$, the observed predictive SD
was $104$ and the predictive median was dragged far negative --- for a
quantity whose data never leave $14.5 \pm 4.4$.

The forecast is the \emph{only} place in the pipeline where $z$ is drawn from
its model distribution rather than read from the stuck chain. It is therefore
the only place the inconsistency can become visible. The forecaster was not the
bug; it was the detector.

\section{Verification}

The chain was reproduced bit-for-bit (same seed, same sequential block order),
then rerun with the single argument \code{z.init = 0.01} removed and nothing
else changed.

\begin{center}
\begin{tabular}{lrrrrrr}
\toprule
& \multicolumn{3}{c}{\textbf{with} \code{z.init = 0.01}} & \multicolumn{3}{c}{\textbf{default} \code{z.init}} \\
\cmidrule(lr){2-4}\cmidrule(lr){5-7}
quantity (truth) & blk 1 & blk 2 & blk 3 & blk 1 & blk 2 & blk 3 \\
\midrule
$\bar z$ \hfill($\approx 1.6$)      & $1.42$ & $0.20$ & $0.22$ & $0.77$ & $1.46$ & $\mathbf{1.68}$ \\
$\lambda$ \hfill($+3$)              & $2.91$ & $-55.8$ & $-97.8$ & $5.81$ & $2.90$ & $\mathbf{2.80}$ \\
$\beta_0$ \hfill($10$)              & $10.36$ & $25.19$ & $36.41$ & $10.34$ & $10.45$ & $\mathbf{10.36}$ \\
$\phi_{1z}$ \hfill($0.80$)          & $0.55$ & $0.94$ & $1.11$ & $0.61$ & $0.58$ & $\mathbf{0.62}$ \\
\midrule
predictive SD, lead 15 \hfill($4.38$) & $5.79$ & $\mathbf{41.9}$ & $\mathbf{104.4}$ & $9.89$ & $5.04$ & $\mathbf{4.94}$ \\
\bottomrule
\end{tabular}
\end{center}

Deleting one argument returns $\lambda$ to $+2.80$, $\beta_0$ to $10.36$, $\bar
z$ to $1.68$, and the lead-15 predictive SD to $4.94$ against a marginal of
$4.38$ --- i.e.\ the forecast now converges to climatology, as an AR forecast
must.

\paragraph{Consequence, in the currency of the study.} Every score in the
previous time-block study was computed from this corrupted predictive sample.
Rescoring the setting-4 study (five blocks $\times$ five datasets) before and
after the one-argument fix makes the damage concrete in the scores themselves,
not merely in the predictive SD:

\begin{center}
\begin{tabular}{lrrr}
\toprule
 & \textbf{before} (\code{z.init=0.01}) & \textbf{after} (fixed) & change \\
\midrule
AR(2) CRPS (mean)        & $5.95$   & $2.64$   & $-56\%$ \\
AR(2) Brier $q_{95}$     & $0.104$  & $0.070$  & $-33\%$ \\
i.i.d.\ CRPS (mean)      & $2.33$   & $2.41$   & $+3\%$ \\
i.i.d.\ Brier $q_{95}$   & $0.053$  & $0.059$  & $+11\%$ \\
\midrule
AR(2)$-$i.i.d.\ CRPS gap  & $+3.62$  & $+0.23$  & \\
\quad (clean 21 blocks)   &          & $-0.015$ ($p = 0.78$) & \\
AR(2)$-$i.i.d.\ Brier gap  & $+0.051$ & $+0.011$ & \\
\quad (clean 21 blocks)   &          & $+0.0001$ ($p = 0.97$) & \\
\bottomrule
\end{tabular}
\end{center}

The i.i.d.\ method barely moves --- it never runs the temporal $z$ update the
bug lives in --- while the AR(2) method's CRPS more than halves and its Brier
falls by a third. Once the four residual reflected blocks are excluded
(\S\ref{sec:residual}), the two methods are a statistical tie on both scores.
The headline finding, ``AR(2) forecasts four times worse than i.i.d.,'' was
never a statement about temporal pooling. It was a statement about an initial
value.

\section{The residual: the same ridge, entered a different way}
\label{sec:residual}

The fix removes the catastrophic corner but not the ridge \eqref{eq:ridge}.
The bug forced \emph{every} chain onto the ridge deterministically, through
the $z$ coordinate; with that door shut, a minority of chains still drift onto
it stochastically during burn-in, through the $(\lambda, \beta_0)$
coordinates. On the guarded block-1 study these ``reflected'' fits appear in
about a quarter of cells: $\lambda$ flips to $-4$ to $-15$, $\beta_0$ rises to
$26$--$35$ so that $\beta_0 + \lambda\bar z$ still reproduces the data mean,
$\bar z$ stays correct (so the level-based assertion below \emph{passes}), and
the predictive SD inflates as $\max_h \sd(\hat y_h) \approx 1.5\,|\hat\lambda|$.
This is a genuine non-identifiability of the model, milder than the bug
(predictive SD $6$--$21$, not $104$), and it is documented separately in the
companion note \code{tex/lambda\_phiz\_ridge}. Two practical consequences
belong here, because both were discovered while cleaning up after this bug.

\subsection{A reflected fit corrupts every proper score, not only the joint
ones}

It is tempting to hope that a latent-process pathology only shows up in
multivariate scores (energy, variogram) that credit temporal dependence, so
that a study reporting the univariate Brier score could ignore it. It cannot.
Refitting the five strongest reflected cells bit-for-bit and scoring the
reflected chain against its healthy replacement on all four scores:

\begin{center}
\small
\begin{tabular}{lrrrrr}
\toprule
cell & $\hat\lambda$ (refl.\,$\to$\,heal.) & Brier $q_{95}$ & CRPS & energy & variogram \\
\midrule
s5 d3 m2 & $-15.2 \to +3.3$ & $0.109 \to 0.132$ & $1.94 \to 2.15$ & $9.9 \to 7.2$ & $243 \to 93$ \\
s4 d6 m2 & $-8.9 \to +2.9$  & $0.138 \to 0.135$ & $5.59 \to 5.48$ & $21.5 \to 18.3$ & $430 \to 302$ \\
s4 d5 m2 & $-7.3 \to +3.9$  & $\mathbf{0.021 \to 0.002}$ & $2.64 \to 1.57$ & $14.3 \to 11.8$ & $276 \to 179$ \\
s4 d5 m1 & $-7.2 \to +3.5$  & $\mathbf{0.027 \to 0.003}$ & $3.19 \to 1.90$ & $14.6 \to 12.2$ & $301 \to 178$ \\
s4 d10 m1 & $-5.0 \to +3.0$ & $0.056 \to 0.023$ & $2.87 \to 2.21$ & $11.6 \to 10.8$ & $234 \to 144$ \\
\midrule
\multicolumn{2}{l}{mean relative damage} & (see text) & $+32\%$ & $+21\%$ & $+78\%$ \\
\bottomrule
\end{tabular}
\end{center}

Every score is worse under reflection --- the inflated $|\lambda|$ widens the
\emph{marginal} forecast at every lead, and no proper score rewards
overdispersion. The ordering, though, is the opposite of the intuition above:
the energy score is the \emph{most} robust ($+21\%$), the pure-dependence
variogram the least ($+78\%$), and CRPS sits between ($+32\%$). The Brier
score is damaged too, and most sharply exactly where it matters: on the three
easy datasets (s4 d5, d10) whose healthy fit gets the tail nearly right, the
reflected chain multiplies the Brier error several-fold ($0.002 \to 0.021$,
$0.003 \to 0.027$, $0.023 \to 0.056$). On the two hard datasets (s5 d3, s4 d6),
where even the healthy fit scores poorly, an overdispersed forecast happens to
hedge, so the Brier barely moves or even improves --- which is not a reprieve
but a warning: a broken fit can win a single validation window by luck, so the
score itself cannot be used to detect the pathology. \emph{Focusing on the
Brier score does not make the identification problem ignorable; it only makes
it intermittent and harder to see.}

\subsection{The same default still lurks in a sibling backend}

The bug's root cause --- a signature default that maps to
$\operatorname{hn.cop}(0,\sigma) = -\infty$ --- was fixed in \code{ar2} but
survives as \code{z.init = 0} in the \code{prop} backend
(\code{code/R/prop/mcmc\_prop.R}). It is currently harmless: \code{prop}
refuses temporal-$z$ blocks with a \code{stop()} guard, and its non-temporal
Gibbs update overwrites the initial value on the first sweep. But it is a
dormant landmine. Enabling temporal $z$ there --- the obvious first step for a
future extension --- would re-arm the original freeze exactly; a direct test
(same shared \code{updateZTS\_AR2} machinery, \code{z.init = 0}) reproduces it
on contact: $z$ frozen, $\phi_z$ variance $0$, \code{NaN} warnings by the
hundred. We have synced the default to \code{NULL} $\mapsto$ the half-normal
median, so the trap is disarmed before anyone steps on it.

\paragraph{Rev.\ 22 July 2026: a third sibling, and this one was live.} The
\code{prop} landmine was dormant behind a guard. The \emph{legacy} backend
(\code{code/R/mcmc\_cont\_lambda.R}), which runs the entire US-all ozone
analysis via \code{package\_load.R}, had no guard: it carried
\code{z.init = 0} in its signature and \emph{does} run the temporal-$z$ block
whenever \code{skew = TRUE} and \code{temporalz = TRUE}. Every legacy
skew-t$+$TS ozone fit to date (settings 51, 52, 54, \ldots, and the MRTS rows
201--208) therefore ran with $z \equiv 0$ frozen for the whole chain and
$\lambda$ wandering its prior --- effectively a symmetric-$t$ fit wearing a
skew-$t$ label. (Predictions were not contaminated, since the skew channel
enters only through $\lambda z_{g} = 0$.) An A/B rerun on a 120-site
subproblem confirms both halves: with \code{z.init = 0} the chain reproduces
the freeze exactly ($z$ all-zero, $\phi_z$ variance $0$, $\lambda$ SD
$19.3$); with the synced default it lands on $\bar z / \E[z] = 0.97$ and
$\lambda$ SD $0.12$. The default is now synced to \code{NULL} $\mapsto$ the
half-normal median, and legacy skew-t$+$TS results produced before and after
this date must not be compared as like for like.

\section{The real lesson}

The original bug was \textbf{loud}: $-\infty$, then \code{NaN}, then $\phi_z$
identically zero for every iteration. It announced itself. Anyone looking at the
output would have stopped.

The workaround silenced it. With \code{z.init = 0.01} there is no $-\infty$, no
\code{NaN}, no all-zero coefficient. There is a chain that mixes, a posterior
that concentrates, a fitted mean that matches the data, and a set of scores that
look entirely publishable. The failure did not go away; it went
\emph{underground}, and it stayed there for a full study.

\begin{quote}
A workaround that removes the \emph{symptom} without removing the \emph{cause}
does not make a system safer. It makes the next failure quieter.
\end{quote}

\paragraph{What should have caught it.} Two assertions, each one line, each
checkable without any knowledge of this bug:

\begin{enumerate}\itemsep2pt
  \item \textbf{Self-consistency of the fit.} The model asserts $z \sim
        \hn(0,\tau^{-1/2})$, hence $\E[z] = \tau^{-1/2}\sqrt{2/\pi}$. The fitted
        $\bar z$ was $7\times$ smaller. \emph{A fit that violates its own
        generative assumption by a factor of seven is not a fit.}
  \item \textbf{Recovery of known truth.} This is a \emph{simulation} study: the
        truth is $\lambda = +3$. The fit returned $\lambda = -98$. The sign alone
        is disqualifying. Simulation studies have ground truth available for free
        and should assert against it automatically.
\end{enumerate}

Neither assertion requires suspecting the initialisation. Either would have
failed on the first block, in minutes, instead of after forty hours and a false
conclusion.

\end{document}
